\documentclass[11pt,a4paper]{article}
\usepackage{operationsform}
\begin{document}
\OPStitle{研究者紹介・解説記事原稿}{TPL-DOC-054}{0.2.0}
\OPSsection{1}{基本情報}
\OPSblanklong{案件・文書の識別子}
\OPSblankshort{作成者と連絡先}
\OPSblankshort{作成日と対象版}
\OPSsection{2}{記事本文}
\OPSblankshort{タイトル案}
\OPSblanklong{リード本文}
\par\Needspace{4\baselineskip}\noindent\textbf{本文と見出し}\par\vspace{1mm}
{\fontsize{8.5}{13}\selectfont
\begin{tabularx}{\linewidth}{@{}XXXX@{}}
\textbf{箇所・ID} & \textbf{記述・出典} & \textbf{確認者・状態} & \textbf{留意点}\\[2mm]
\strut & \strut & \strut & \strut\\[5mm]
\strut & \strut & \strut & \strut\\[5mm]
\strut & \strut & \strut & \strut\\[5mm]
\end{tabularx}
}\vspace{2mm}\par
\OPSblanklong{引用・図表の挿入位置}
\OPSsection{3}{根拠と編集記録}
\OPSblanklong{主要主張と資料参照}
\OPSblanklong{未確定表現と留保}
\OPSblanklong{編集者への確認事項}
\OPSblanklong{改稿履歴}
\OPSsection{4}{確認と履歴}
\par\Needspace{4\baselineskip}\noindent\textbf{参照資料と根拠ID}\par\vspace{1mm}
{\fontsize{8.5}{13}\selectfont
\begin{tabularx}{\linewidth}{@{}XXXX@{}}
\textbf{箇所・ID} & \textbf{記述・出典} & \textbf{確認者・状態} & \textbf{留意点}\\[2mm]
\strut & \strut & \strut & \strut\\[5mm]
\strut & \strut & \strut & \strut\\[5mm]
\strut & \strut & \strut & \strut\\[5mm]
\end{tabularx}
}\vspace{2mm}\par
\OPSblanklong{未確認事項・例外}
\par\Needspace{4\baselineskip}\noindent\textbf{点検・改訂の記録}\par\vspace{1mm}
{\fontsize{8.5}{13}\selectfont
\begin{tabularx}{\linewidth}{@{}XXXX@{}}
\textbf{箇所・ID} & \textbf{記述・出典} & \textbf{確認者・状態} & \textbf{留意点}\\[2mm]
\strut & \strut & \strut & \strut\\[5mm]
\strut & \strut & \strut & \strut\\[5mm]
\strut & \strut & \strut & \strut\\[5mm]
\end{tabularx}
}\vspace{2mm}\par
\end{document}
